SemEval-2026 Task~6 (CLARITY) asks how a politician's answer responds to an interview question.
Subtask~2 has nine evasion types, and each type fixes one of the three clarity levels of Subtask~1.
We set out to find what limits single-pass classifiers on this task, changing one thing at a time and writing down our predictions before every run.
A DeBERTa-v3-large classifier built this way reaches 0.384 multi-reference macro-F1 on Subtask~2 per model, and 0.405 as a ten-seed system with a one-parameter logit adjustment.
Re-ranking, three kinds of hierarchy, rebalanced losses and model soups all failed to improve on it, and the error analysis points at the model rather than at label noise.
When we swap only the backbone for Qwen3-8B fine-tuned with LoRA, keeping the data, input, loss and model selection fixed, Subtask~2 rises by 0.092 and Subtask~1 by 0.095 on every one of ten paired seeds, and the ten-seed system reaches 0.543 and 0.746 (+0.136 on Subtask~2, 95\% CI [0.054, 0.224]).
The gain is spread over most classes instead of sitting on the label boundaries we had expected.
All numbers are on the development set, because the test labels were never released.
